\documentclass[border=8pt]{standalone}
\usepackage{fontspec}
\defaultfontfeatures{Ligatures=TeX}
\setmainfont{Palatino Linotype}
\setsansfont{Arial}
\usepackage{amsmath,amssymb}
% =====================================================================
% CAS PCBA Engineer — figures preamble (TikZ schematics)
% Little Henderson · Noetic Compass Press · 2026
%
% PREAMBLE-SAFE: this file is \input into the book preamble by the build.
% It must contain ONLY preamble code (package loads, color defs, \tikzset).
% No \begin{document}, no document-body material.
%
% Visual grammar (consistent across all 7 figures):
%   Domain colors as light fills + dark borders + black text:
%     P (Principled) = gold/amber, C (Cognitive) = blue,
%     E (Evaluative) = red,        O (Operational) = green.
%   Neutral / non-domain blocks = gray.
%   Rounded-rectangle blocks; clean downward signal-flow arrows (-{Latex}).
%   Designed to stay <= ~11cm wide (6x9 measure) and legible in grayscale
%   (rely on borders + labels, not color alone).
% =====================================================================

\usepackage{tikz}
\usetikzlibrary{arrows.meta,positioning,shapes.geometric,calc,fit,backgrounds}

% --- pgfplots for the vector charts (CJT curves, Bode, filters, decay) --
\usepackage{pgfplots}
\pgfplotsset{compat=1.18}
\usepgfplotslibrary{groupplots}            % stacked Bode panels share an x-axis
\pgfplotsset{
  % Shared house style for the analytic charts: light grid, sans labels,
  % measure-friendly default size, grayscale-legible line cycle.
  casaxis/.style={
    width=10.6cm, height=6.2cm,
    font=\sffamily\footnotesize,
    axis line style={draw=casNedge, line width=0.5pt},
    tick label style={font=\sffamily\scriptsize},
    label style={font=\sffamily\footnotesize},
    title style={font=\sffamily\footnotesize\bfseries},
    every axis grid/.style={draw=casNedge!25, line width=0.3pt},
    grid=both, tick align=outside, tickpos=left,
    legend cell align=left,
    legend style={font=\sffamily\scriptsize, draw=casNedge, fill=white,
      inner sep=2pt, row sep=0.5pt}
  }
}

% --- Domain palette (light fills, dark borders) ------------------------
\definecolor{casPfill}{RGB}{255,243,205}   % amber tint  (Principled)
\definecolor{casPedge}{RGB}{153,110,0}     % dark amber
\definecolor{casCfill}{RGB}{213,230,247}   % blue tint   (Cognitive)
\definecolor{casCedge}{RGB}{26,76,140}     % dark blue
\definecolor{casEfill}{RGB}{250,219,216}   % red tint    (Evaluative)
\definecolor{casEedge}{RGB}{150,40,30}     % dark red
\definecolor{casOfill}{RGB}{215,239,221}   % green tint  (Operational)
\definecolor{casOedge}{RGB}{30,110,55}     % dark green
\definecolor{casNfill}{RGB}{233,233,233}   % gray tint   (neutral)
\definecolor{casNedge}{RGB}{90,90,90}      % dark gray

% Accent for "go / pass" and "stop / refuse" paths (still labelled in text,
% so grayscale readers never depend on the color).
\definecolor{casGreen}{RGB}{30,110,55}
\definecolor{casRed}{RGB}{150,40,30}

% --- Reusable block + flow styles --------------------------------------
\tikzset{
  % Base block: rounded rectangle, black text, fixed measure-friendly width.
  casblock/.style={
    rectangle, rounded corners=2pt, draw, line width=0.6pt,
    text=black, align=center, inner sep=3pt,
    minimum height=8mm, font=\sffamily\footnotesize,
    fill=casNfill, draw=casNedge
  },
  % Domain-tinted blocks (inherit casblock, override fill+border).
  pblock/.style={casblock, fill=casPfill, draw=casPedge},
  cblock/.style={casblock, fill=casCfill, draw=casCedge},
  eblock/.style={casblock, fill=casEfill, draw=casEedge},
  oblock/.style={casblock, fill=casOfill, draw=casOedge},
  nblock/.style={casblock, fill=casNfill, draw=casNedge},
  % Wider container block for multi-line stage descriptions.
  widecasblock/.style={casblock, text width=58mm},
  widepblock/.style={pblock, text width=58mm},
  widecblock/.style={cblock, text width=58mm},
  wideeblock/.style={eblock, text width=58mm},
  wideoblock/.style={oblock, text width=58mm},
  widenblock/.style={nblock, text width=58mm},
  % Terminal nodes (signal entry / exit): small pill.
  casterm/.style={casblock, fill=white, draw=black, font=\sffamily\scriptsize\bfseries,
    minimum height=6mm, rounded corners=3pt},
  % Signal-flow arrow (downward by default).
  flow/.style={-{Latex[length=2mm,width=1.6mm]}, line width=0.6pt, draw=black},
  % Feedback / secondary path (dashed).
  fbflow/.style={-{Latex[length=2mm,width=1.6mm]}, line width=0.5pt, draw=casNedge,
    dashed},
  % Small italic edge label.
  flowlbl/.style={font=\sffamily\scriptsize\itshape, fill=white, inner sep=1pt},
  % Group / section frame drawn behind blocks (backgrounds layer).
  groupframe/.style={rounded corners=3pt, draw=casNedge, line width=0.5pt,
    dash pattern=on 1.5pt off 1.5pt, fill=casNfill, fill opacity=0.25},
  grouplbl/.style={font=\sffamily\scriptsize\bfseries, text=casNedge}
}

% Helper: render a CAS state literal in monospace with carets escaped, in
% math-free text mode. Usage: \casnote{^8^C|A}  ->  prints ^8^C|A safely.
\providecommand{\casnote}[1]{{\ttfamily\scriptsize #1}}

\begin{document}
\centering
\begin{tikzpicture}
\begin{axis}[casaxis, width=8.8cm, height=6.0cm,
    xlabel={cause activation}, ylabel={effect activation},
    xmin=0, xmax=1.0, ymin=0, ymax=1.0,
    xtick={0,0.2,0.4,0.6,0.8,1.0}, ytick={0,0.2,0.4,0.6,0.8,1.0},
    legend style={at={(1.02,1.0)}, anchor=north west}, legend columns=1,
    clip mode=individual,
  ]
  % four bias curves: cutoff at 0.1, linear active (slope rises w/ bias),
  % saturation knee (plateau 0.30/0.55/0.78/0.95), slight early-effect slope.
  \addplot[black, densely dotted, line width=1pt, domain=0.1:1, samples=120]
    {(x<0.6455)*(0.55*(x-0.1)) + (x>=0.6455)*(0.30 + 0.05*(x-0.6455))};
  \addplot[casCedge, dashed, line width=1pt, domain=0.1:1, samples=120]
    {(x<0.6789)*(0.95*(x-0.1)) + (x>=0.6789)*(0.55 + 0.06*(x-0.6789))};
  \addplot[casOedge, dashdotted, line width=1.1pt, domain=0.1:1, samples=120]
    {(x<0.6778)*(1.35*(x-0.1)) + (x>=0.6778)*(0.78 + 0.07*(x-0.6778))};
  \addplot[casEedge, solid, line width=1.1pt, domain=0.1:1, samples=120]
    {(x<0.6429)*(1.75*(x-0.1)) + (x>=0.6429)*(0.95 + 0.04*(x-0.6429))};
  \legend{bias=0.2, bias=0.4, bias=0.6, bias=0.8}
  % load line (R=30), crosses the bias=0.6 curve at the Q-point (~0.43, ~0.45)
  \addplot[casNedge, line width=0.9pt, domain=0:0.93, samples=2] {0.837 - 0.9*x};
  \node[casNedge, font=\sffamily\scriptsize, anchor=north west] at (axis cs:0.6,0.30) {load line ($R{=}30$)};
  \addplot[only marks, mark=*, mark size=2pt, black] coordinates {(0.43,0.45)};
  \node[font=\sffamily\scriptsize\bfseries, anchor=south east] at (axis cs:0.42,0.46) {Q};
  \node[font=\sffamily\scriptsize, anchor=north east] at (axis cs:0.42,0.44) {(0.43, 0.45)};
  \node[font=\sffamily\scriptsize\itshape, anchor=north] at (axis cs:0.1,0.12) {cutoff};
  \draw[casNedge, line width=0.4pt, densely dotted] (axis cs:0.1,0) -- (axis cs:0.1,1.0);
  \node[font=\sffamily\scriptsize\itshape, anchor=south] at (axis cs:0.2,0.78) {active region};
  \draw[-{Latex[length=1.6mm]}, casNedge, line width=0.4pt]
    (axis cs:0.2,0.76) -- (axis cs:0.27,0.55);
  \node[font=\sffamily\scriptsize\itshape, anchor=south] at (axis cs:0.62,0.99) {saturation knee};
  \draw[-{Latex[length=1.6mm]}, casNedge, line width=0.4pt]
    (axis cs:0.62,0.98) -- (axis cs:0.645,0.83);
\end{axis}
\end{tikzpicture}
\end{document}
